\documentclass[../../main.tex]{subfiles}
\begin{document}

\begin{center}
  {\fontsize{48pt}{58pt}\selectfont \textbf{1977}\normalsize\textbf{年度}}
\end{center}

\begin{table}[htb]
  \centering
  \caption{出題テーマと難易度}
  \begin{tabular}{|c|p{9cm}|c|}
    \hline
    \textbf{問題番号} & \textbf{テーマ} & \textbf{難易度} \\
    \hline
    第1問 & 絶対値を含む三次関数の最大値の最小化 & \\
    \hline
    第2問 & 有理数座標の二点がなす角 & \\
    \hline
    第3問 & 半平面に含まれる円の最大半径 & \\
    \hline
    第4問 & 転がる正方形の点の軌跡と回転体 & \\
    \hline
    第5問 & 新課程：行列と円周上の軌跡／旧課程：円に内接する五角形 & \\
    \hline
    第6問 & 新課程：空間の直線と平面／旧課程：三角関数の極値と変曲点 & \\
    \hline
  \end{tabular}
\end{table}

\begin{table}[htb]
  \centering
  \caption{解答}
  \renewcommand{\arraystretch}{1.4}
  \begin{tabular}{|c|c|p{9cm}|}
    \hline
    \textbf{問題番号} & \textbf{小問} & \textbf{答え} \\
    \hline
    第1問 & -- & $k=\dfrac{1}{4}$で最小値$\dfrac{1}{4}$ \\
    \hline
    第2問 & -- & $ab$，$\cos\theta$，$\sin\theta$の有理性は同値（証明問題） \\
    \hline
    第3問 & (1) & $r(P)\ge 2\sqrt{3}-3$ \\
    \cline{2-3}
    & (2) & $r(P)\ge\dfrac{\sqrt{3}-1}{2}$ \\
    \hline
    第4問 & -- & $V(a)=\pi\left\{12(a-1)^2+\dfrac{40}{3}\right\}$．$a=1$で最小値$\dfrac{40\pi}{3}$ \\
    \hline
    第5問 & 新課程 & $a=d=\dfrac{1-t^2}{1+t^2}$，$b=-c=-\dfrac{2t}{1+t^2}$．軌跡は円$x^2+y^2=25$から$(-3,-4)$を除いた部分 \\
    \cline{2-3}
    & 旧課程 & 与えられた角の条件が必要十分（証明問題） \\
    \hline
    第6問 & 新課程(1) & 平面は$-2x-y+2z+3=0$ \\
    \cline{2-3}
    & 新課程(2) & $l$との交点は$(3,1,2)$，$m$との交点は$(1,-1,0)$ \\
    \cline{2-3}
    & 旧課程(1) & 最小値は$-1-\dfrac{a}{4}$．最大値は$0<a\le1$で$1-\dfrac{a}{4}$，$a>1$で$\dfrac{1}{2a}+\dfrac{a}{4}$ \\
    \cline{2-3}
    & 旧課程(2) & 変曲点は$0<a\le1$で$2$個，$a>1$で$4$個（凹凸は解答参照） \\
    \hline
  \end{tabular}
\end{table}

\paragraph{試験形態}
6問．第5・6問は新旧課程別．試験時間は要確認．

\end{document}
